% ============================================================
% SRT v10 Updates — changes from implementation experience
% Apply these edits to srt_v9.tex
% ============================================================

% ============================================================
% UPDATE 1: On-disk format (Appendix A, ~line 923)
% Add directory sibling records to the storage description
% ============================================================
% REPLACE lines 925-930:

Summary records are stored inside the repository in colocated hidden directories:
\begin{itemize}[nosep]
    \item Each directory in the repository contains a hidden \texttt{.srt/} folder.
    \item The directory's own summary record is stored as \texttt{.srt/\_\_dir\_\_.md}.
    \item Each direct file leaf in that directory is stored as \texttt{.srt/<filename>.md}.
    \item Each direct subdirectory child is \emph{also} stored as a sibling record at the parent level: \texttt{.srt/<dirname>.md}. This sibling duplicates the content of \texttt{<dirname>/.srt/\_\_dir\_\_.md} but lives alongside file records so that all immediate children---files and directories alike---are enumerable from a single \texttt{.srt/} folder. This is necessary for embedding-assisted routing (\Cref{app:embeddings}), where each child must have a colocated vector sidecar.
\end{itemize}

% ============================================================
% UPDATE 2: Traversal (Section 4.2, ~line 309)
% Add embedding-based traversal as an alternative policy
% ============================================================
% AFTER the existing \paragraph{Traversal policy.} block (line 309-314),
% ADD a new paragraph:

\paragraph{Embedding-based traversal.}
An alternative traversal policy replaces the per-level LLM oracle with precomputed vector similarity.
During an offline embedding pass, each node's summary is encoded into a dense vector stored as a sidecar file alongside the summary record (\Cref{app:embeddings}).
At query time, the query is embedded once, and at each internal node the system ranks children by cosine similarity between the query vector and each child's precomputed vector.
The top-$k$ children (a beam width parameter) are selected for descent without an LLM call.

This policy is a complete traversal strategy, not only a pre-filter.
In the reference implementation, a single CLI invocation performs beam-search descent from root to leaves in under 100\,ms on a 300-node tree, compared to seconds per level for the LLM-oracle policy.
The tradeoff is that cosine similarity over summary embeddings is a weaker relevance signal than a full LLM judgment; the embedding policy may miss children whose relevance depends on reasoning not captured by vector proximity.
A hybrid policy---embedding pre-filter followed by LLM confirmation on the survivors---combines the speed of embedding routing with the judgment of the oracle.

% ============================================================
% UPDATE 3: Appendix B title and framing (~line 1069)
% Reframe from "optimization" to "implementation" since it
% works as standalone routing, not just pre-filtering
% ============================================================
% REPLACE line 1069:

\section{Embedding-Assisted Routing}\label{app:embeddings}

% REPLACE lines 1071 (the opening paragraph):

The reference design routes queries through LLM calls: at each tree level, the agent reads child summaries and asks the LLM which branches are relevant.
This appendix describes an embedding layer that accelerates or replaces that per-level LLM call.
Precomputed vector embeddings of node summaries enable cosine-similarity ranking of children, which can serve as either a pre-filter before the LLM oracle or as a standalone traversal policy.

% ============================================================
% UPDATE 4: Appendix B — child selection (~line 1079-1089)
% Update to acknowledge standalone mode
% ============================================================
% REPLACE lines 1087-1089:

Children with similarity below a threshold $\tau$ can be pruned before the LLM call, reducing the number of summaries the LLM must read.
Children above $\tau$ are passed to the LLM for the final selection decision.

When used as a pre-filter, the LLM still makes the routing judgment but over a smaller candidate set.
The embedding layer can also operate as a standalone traversal policy by selecting the top-$k$ children at each level without invoking the LLM at all.
In practice, beam-search descent with $k = 3$ and no LLM calls identifies the correct target files for a range of structural queries, completing full root-to-leaf traversal in under 100\,ms.
The LLM oracle remains available for queries where vector similarity is insufficient---for example, queries requiring multi-hop reasoning about relationships between distant modules.

% ============================================================
% UPDATE 5: Appendix B — Storage (~line 1096-1099)
% Pin the .vec sidecar format
% ============================================================
% REPLACE lines 1096-1099:

\subsection{Storage}

The reference implementation stores embeddings as per-node JSON sidecar files colocated with summary records.
For a summary record at \texttt{.srt/foo.py.md}, the corresponding embedding is stored at \texttt{.srt/foo.py.vec} containing:

\begin{verbatim}
{
  "model": "BAAI/bge-small-en-v1.5",
  "content_hash": "e5f6a7b8c9d0...",
  "vector": [0.0123, -0.0456, ...]
}
\end{verbatim}

The \texttt{content\_hash} field mirrors the summary record's hash; a mismatch indicates the embedding is stale and must be recomputed.
The \texttt{model} field records which embedding model produced the vector, enabling detection when the model changes (requiring a full re-embedding pass).
Incremental updates follow the same hash-based invalidation as summaries: when a node's content hash changes, its embedding is recomputed.

This sidecar design preserves the git-native storage property: embedding files are plain text tracked alongside summaries, visible in diffs, and require no external infrastructure.
The tradeoff relative to a SQLite store is a larger number of small files; in practice the per-file overhead is negligible and the gain in diff visibility and per-node invalidation granularity is worth it.

% ============================================================
% UPDATE 6: Appendix A — Design scope and limitations (~line 1052)
% Update the limitation about LLM-only routing
% ============================================================
% REPLACE the fourth bullet (line 1059):

    \item The default traversal relies on an LLM to decide which child summaries are worth descending into; an embedding-based alternative is described in \Cref{app:embeddings} and implemented in the reference indexer.

% REPLACE lines 1062-1063:

These are implementation choices, not limitations of the \srt{} abstraction itself.
A richer implementation could add symbol-level leaves, structured summary schemas, or batched summarization for wide directories without changing the core data structure or query protocol.

% ============================================================
% UPDATE 7: Appendix A — Reference implementation CLI (~line 1017)
% Add embed/query/route commands to the deployment section
% ============================================================
% AFTER the existing query-time protocol enumeration (line 1039),
% ADD a new paragraph:

\paragraph{CLI tools for embedding-assisted routing.}
The reference indexer provides three additional commands that support the embedding layer:
\begin{itemize}[nosep]
    \item \texttt{semtree embed [path]} computes embeddings for all existing summary records, writing \texttt{.vec} sidecar files. Skips nodes whose embedding is already fresh (matching content hash and model).
    \item \texttt{semtree query <query> [path]} ranks the immediate children of a directory by cosine similarity to a natural-language query. Useful for one-level pre-filtering.
    \item \texttt{semtree route <query> [path]} performs a complete beam-search descent from root to leaves in a single invocation, ranking children by cosine similarity at each level. Returns the traversal path and candidate files.
\end{itemize}
The \texttt{route} command is the primary entry point for embedding-based traversal: it loads the embedding model once, descends through the tree, and reports candidates in under 100\,ms for trees of several hundred nodes.
The \texttt{build} command optionally runs the embedding pass after summarization (controlled by a \texttt{--no-embed} flag).
